La descripción del movimiento realizada en el capítulo anterior nos permitió
establecer un sistema de referencia, definir la posición de una partícula y
distinguir entre distancia recorrida y desplazamiento.

Ahora queremos avanzar un paso más.

No solamente nos interesa saber dónde se encuentra un cuerpo o cuánto cambió
su posición. También queremos saber \textbf{qué tan rápidamente cambia esa
posición con el tiempo}.

Esta pregunta conduce al concepto de velocidad.

En este capítulo estudiaremos movimientos cuya trayectoria es una línea
recta. Aunque se trata del caso geométricamente más sencillo, aquí aparecen
conceptos fundamentales que utilizaremos posteriormente en prácticamente
toda la mecánica.


\section{Movimiento en una dimensión}

Cuando una partícula se mueve sobre una trayectoria recta podemos hacer
coincidir uno de los ejes de coordenadas con esa trayectoria.

Por ejemplo, si el movimiento ocurre horizontalmente, podemos utilizar el
eje \(x\):

\begin{center}
\begin{tikzpicture}
    \draw[-{Latex},thick] (-5,0) -- (5,0) node[right] {$x$};

    \foreach \x in {-4,-2,0,2,4}
    {
        \draw (\x,-0.12) -- (\x,0.12);
    }

    \node[below] at (0,-0.15) {$0$};
    \node[above] at (4,0.15) {sentido positivo};
\end{tikzpicture}
\end{center}

Una vez elegido el sentido positivo, cualquier posición puede representarse
mediante una sola coordenada \(x\).

La posición de la partícula puede cambiar con el tiempo, por lo que escribimos

\[ x=x(t). \]

\begin{interpretacion}
La expresión \(x(t)\) contiene la información fundamental del movimiento
rectilíneo.

Si conocemos \(x(t)\), sabemos dónde se encuentra la partícula en cada
instante y, a partir de esa información, podemos estudiar su velocidad y
posteriormente su aceleración.
\end{interpretacion}


\section{Velocidad media}

Supongamos que una partícula se encuentra inicialmente en la posición
\(x_i\) en el instante \(t_i\) y posteriormente se encuentra en la posición
\(x_f\) en el instante \(t_f\).

El desplazamiento es

\[ \Delta x=x_f-x_i, \]

y el intervalo de tiempo es

\[ \Delta t=t_f-t_i. \]

Queremos comparar cuánto cambió la posición con el tiempo empleado en
producir ese cambio.

\begin{definicion}[Velocidad media]
La \textbf{velocidad media} de una partícula durante un intervalo de tiempo
se define como el cociente entre su desplazamiento y el intervalo de tiempo
transcurrido.
\end{definicion}

En movimiento rectilíneo,

\[ \boxed{\bar{v}_x=\frac{\Delta x}{\Delta t}
=\frac{x_f-x_i}{t_f-t_i}}. \]

Aquí:

\begin{itemize}
    \item \(\bar{v}_x\) es la velocidad media;
    \item \(\Delta x\) es el desplazamiento;
    \item \(\Delta t\) es el intervalo de tiempo.
\end{itemize}

En el Sistema Internacional,

\[ [\bar{v}_x]=\si{\meter\per\second}. \]

\begin{interpretacion}
La velocidad media indica el cambio de posición que correspondería, en
promedio, a cada unidad de tiempo durante el intervalo considerado.

No nos dice necesariamente cómo se movió el cuerpo en cada instante.
\end{interpretacion}


\subsection{El signo de la velocidad media}

Como siempre consideramos

\[ \Delta t>0, \]

el signo de la velocidad media depende del signo del desplazamiento.

Si

\[ \Delta x>0, \]

entonces

\[ \bar{v}_x>0. \]

El desplazamiento neto ocurrió en el sentido positivo del eje.

Si

\[ \Delta x<0, \]

entonces

\[ \bar{v}_x<0. \]

El desplazamiento neto ocurrió en el sentido negativo.

Si

\[ \Delta x=0, \]

entonces

\[ \bar{v}_x=0. \]

\begin{advertencia}
Una velocidad media igual a cero no significa necesariamente que el cuerpo
permaneció en reposo.

Puede haberse movido y posteriormente haber regresado a su posición inicial.
\end{advertencia}


\begin{ejemplo}[Velocidad media positiva]

Un automóvil se encuentra en

\[ x_i=-\SI{20}{\meter} \]

cuando

\[ t_i=\SI{2}{\second}, \]

y posteriormente se encuentra en

\[ x_f=\SI{40}{\meter} \]

cuando

\[ t_f=\SI{7}{\second}. \]

El desplazamiento es

\[
\Delta x=x_f-x_i
=\SI{40}{\meter}-(-\SI{20}{\meter})
=\SI{60}{\meter}.
\]

El intervalo de tiempo es

\[
\Delta t=t_f-t_i
=\SI{7}{\second}-\SI{2}{\second}
=\SI{5}{\second}.
\]

Por tanto,

\[
\bar{v}_x
=\frac{\Delta x}{\Delta t}
=\frac{\SI{60}{\meter}}{\SI{5}{\second}}
=\SI{12}{\meter\per\second}.
\]

Así,

\[ \boxed{\bar{v}_x=+\SI{12}{\meter\per\second}}. \]

El signo positivo indica que el desplazamiento neto ocurrió en el sentido
positivo del eje.
\end{ejemplo}


\begin{ejemplo}[Velocidad media negativa]

Una motocicleta se encuentra inicialmente en

\[ x_i=\SI{50}{\meter} \]

y, \(\SI{10}{\second}\) después, se encuentra en

\[ x_f=-\SI{30}{\meter}. \]

Su desplazamiento es

\[
\Delta x
=x_f-x_i
=-\SI{30}{\meter}-\SI{50}{\meter}
=-\SI{80}{\meter}.
\]

Por tanto,

\[
\bar{v}_x
=\frac{-\SI{80}{\meter}}{\SI{10}{\second}}
=-\SI{8.0}{\meter\per\second}.
\]

Entonces,

\[ \boxed{\bar{v}_x=-\SI{8.0}{\meter\per\second}}. \]

El resultado negativo no significa que la motocicleta se moviera
``lentamente''.

El signo informa sobre el sentido del desplazamiento neto.
\end{ejemplo}


\section{Rapidez media}

En el lenguaje cotidiano solemos utilizar rapidez y velocidad como si fueran
sinónimos.

En Física debemos distinguirlas.

La velocidad media se calcula utilizando el \textbf{desplazamiento}.

La rapidez media se calcula utilizando la \textbf{distancia total recorrida}.

\begin{definicion}[Rapidez media]
La \textbf{rapidez media} es el cociente entre la distancia total recorrida
y el intervalo de tiempo empleado.
\end{definicion}

Si \(s\) representa la distancia total recorrida,

\[ \boxed{\bar{v}_{\text{rap}}=\frac{s}{\Delta t}}. \]

La rapidez media es una magnitud escalar y, por tanto,

\[ \bar{v}_{\text{rap}}\geq 0. \]


\subsection{Velocidad media y rapidez media}

Las expresiones

\[ \bar{v}_x=\frac{\Delta x}{\Delta t} \]

y

\[ \bar{v}_{\text{rap}}=\frac{s}{\Delta t} \]

parecen similares, pero describen cantidades diferentes.

La diferencia fundamental está en el numerador:

\[
\text{velocidad media}
\quad\longrightarrow\quad
\text{desplazamiento},
\]

mientras que

\[
\text{rapidez media}
\quad\longrightarrow\quad
\text{distancia recorrida}.
\]

\begin{idea}
En general,

\[ \boxed{\bar{v}_{\text{rap}}\neq\abs{\bar{v}_x}}. \]

Solo coinciden cuando la distancia recorrida es igual a la magnitud del
desplazamiento, como ocurre en un movimiento rectilíneo sin cambio de
sentido durante el intervalo considerado.
\end{idea}


\begin{ejemplo}[Rapidez media y velocidad media diferentes]

Una persona camina

\[ \SI{60}{\meter} \]

hacia la derecha y después

\[ \SI{20}{\meter} \]

hacia la izquierda.

El recorrido completo dura

\[ \SI{40}{\second}. \]

Elegimos la derecha como sentido positivo.

La distancia total es

\[
s=\SI{60}{\meter}+\SI{20}{\meter}
=\SI{80}{\meter}.
\]

La rapidez media es

\[
\bar{v}_{\text{rap}}
=\frac{\SI{80}{\meter}}{\SI{40}{\second}}
=\SI{2.0}{\meter\per\second}.
\]

El desplazamiento es

\[
\Delta x
=\SI{60}{\meter}-\SI{20}{\meter}
=\SI{40}{\meter}.
\]

Por tanto,

\[
\bar{v}_x
=\frac{\SI{40}{\meter}}{\SI{40}{\second}}
=\SI{1.0}{\meter\per\second}.
\]

Así,

\[
\boxed{\bar{v}_{\text{rap}}
=\SI{2.0}{\meter\per\second}}
\]

y

\[
\boxed{\bar{v}_x
=+\SI{1.0}{\meter\per\second}}.
\]

La diferencia aparece porque la distancia recorrida y el desplazamiento no
son iguales.
\end{ejemplo}


\begin{ejemplo}[Velocidad media cero con rapidez media distinta de cero]

Un corredor avanza

\[ \SI{100}{\meter} \]

y posteriormente regresa hasta el punto de partida.

El recorrido completo dura

\[ \SI{50}{\second}. \]

La distancia recorrida es

\[
s=\SI{200}{\meter}.
\]

Por tanto,

\[
\bar{v}_{\text{rap}}
=\frac{\SI{200}{\meter}}{\SI{50}{\second}}
=\SI{4.0}{\meter\per\second}.
\]

Como termina en el mismo punto donde comenzó,

\[ \Delta x=0. \]

Entonces,

\[
\bar{v}_x=0.
\]

Por tanto,

\[
\boxed{\bar{v}_{\text{rap}}
=\SI{4.0}{\meter\per\second}},
\qquad
\boxed{\bar{v}_x=0}.
\]

Este ejemplo muestra por qué la rapidez media no debe calcularse tomando
simplemente el valor absoluto de la velocidad media.
\end{ejemplo}


\section{Velocidad instantánea}

La velocidad media describe un intervalo completo.

Pero imaginemos que un automóvil realiza un viaje de una hora.

Saber su velocidad media durante toda la hora no nos dice necesariamente qué
velocidad tenía exactamente a los veinte minutos.

Para describir el movimiento en un instante determinado necesitamos el
concepto de \textbf{velocidad instantánea}.

\begin{notapedagogica}
Puede introducirse esta idea preguntando:

\begin{quote}
¿Qué indica realmente el velocímetro de un automóvil?
\end{quote}

No está mostrando la velocidad media desde que comenzó el viaje.

Está indicando qué tan rápido se mueve el automóvil en ese momento.
\end{notapedagogica}


\subsection{De un intervalo a un instante}

Supongamos que queremos conocer la velocidad de una partícula en el instante
\(t\).

Podemos calcular primero la velocidad media en un intervalo relativamente
grande alrededor de ese instante.

Después reducimos el intervalo.

Por ejemplo:

\[
\Delta t=\SI{1}{\second},
\]

luego

\[
\Delta t=\SI{0.1}{\second},
\]

después

\[
\Delta t=\SI{0.01}{\second},
\]

y así sucesivamente.

A medida que el intervalo se hace más pequeño, la velocidad media describe
cada vez mejor el movimiento alrededor del instante considerado.

\begin{definicion}[Velocidad instantánea]
La \textbf{velocidad instantánea} es el límite de la velocidad media cuando
el intervalo de tiempo tiende a cero.
\end{definicion}

Matemáticamente,

\[ \boxed{v_x=\lim_{\Delta t\to0}\frac{\Delta x}{\Delta t}}. \]

En lenguaje del cálculo diferencial,

\[ \boxed{v_x=\frac{\dd x}{\dd t}}. \]

\begin{interpretacion}
La velocidad instantánea es la tasa instantánea de cambio de la posición
respecto del tiempo.

Nos indica qué tan rápidamente está cambiando la posición y en qué sentido
lo hace en un instante determinado.
\end{interpretacion}

\begin{nota}
La definición mediante derivadas permite describir con precisión movimientos
en los que la velocidad cambia continuamente.

Cuando todavía no se dispone de cálculo diferencial, la misma idea puede
comprenderse y utilizarse gráficamente mediante la pendiente de una recta
tangente a la gráfica posición-tiempo.
\end{nota}


\subsection{Rapidez instantánea}

La rapidez instantánea es la magnitud de la velocidad instantánea.

En una dimensión,

\[ \boxed{v=\abs{v_x}}. \]

Por tanto,

\[ v\geq0. \]

Por ejemplo, las velocidades

\[ v_x=+\SI{15}{\meter\per\second} \]

y

\[ v_x=-\SI{15}{\meter\per\second} \]

corresponden a movimientos en sentidos opuestos, pero ambos tienen la misma
rapidez:

\[ v=\SI{15}{\meter\per\second}. \]


\begin{ejemplo}[Velocidad instantánea a partir de una función]

La posición de una partícula está dada por

\[ x(t)=2t^2+3t-1, \]

donde \(x\) se expresa en metros y \(t\) en segundos.

Queremos determinar su velocidad instantánea.

Por definición,

\[
v_x=\frac{\dd x}{\dd t}.
\]

Derivando,

\[
v_x(t)
=\frac{\dd}{\dd t}\left(2t^2+3t-1\right),
\]

por lo que

\[ v_x(t)=4t+3. \]

Así,

\[ \boxed{v_x(t)=4t+3}. \]

En \(t=\SI{2}{\second}\),

\[
v_x(2)
=4(2)+3
=11.
\]

Por tanto,

\[ \boxed{v_x=\SI{11}{\meter\per\second}}. \]

El signo positivo indica que en ese instante la posición está aumentando y
la partícula se mueve en el sentido positivo del eje.
\end{ejemplo}


\section{Interpretación de una gráfica posición-tiempo}

Una de las formas más útiles de representar el movimiento consiste en
construir una gráfica de posición \(x\) en función del tiempo \(t\).

En el eje horizontal colocamos el tiempo.

En el eje vertical colocamos la posición.

\begin{figure}[H]
    \centering

    \begin{tikzpicture}
        \begin{axis}[
            axis lines=middle,
            xlabel={$t$},
            ylabel={$x$},
            xmin=0,
            xmax=6,
            ymin=-1,
            ymax=8,
            width=10cm,
            height=7cm,
            xtick={0,1,2,3,4,5},
            ytick={0,2,4,6}
        ]

            \addplot[
                domain=0:5,
                samples=2,
                thick
            ] {1+1.2*x};

        \end{axis}
    \end{tikzpicture}

    \caption{Ejemplo de gráfica de posición en función del tiempo.}
    \label{fig:posicion_tiempo_rectilineo}
\end{figure}

Cada punto de la gráfica responde una pregunta:

\begin{quote}
¿Dónde se encontraba la partícula en determinado instante?
\end{quote}

Sin embargo, la gráfica contiene aún más información.

Su pendiente está relacionada con la velocidad.


\subsection{Pendiente y velocidad media}

Consideremos dos puntos de la gráfica:

\[
(t_i,x_i)
\]

y

\[
(t_f,x_f).
\]

La pendiente de la recta secante que une estos puntos es

\[
m
=
\frac{x_f-x_i}{t_f-t_i}.
\]

Pero

\[
x_f-x_i=\Delta x
\]

y

\[
t_f-t_i=\Delta t.
\]

Por tanto,

\[ \boxed{m=\frac{\Delta x}{\Delta t}=\bar{v}_x}. \]

\begin{idea}
En una gráfica de posición contra tiempo,

\[
\boxed{
\text{pendiente de una secante}
=
\text{velocidad media}
}.
\]
\end{idea}


\subsection{Pendiente y velocidad instantánea}

Si hacemos que los dos puntos de la secante se aproximen cada vez más entre
sí, la recta secante se aproxima a la recta tangente.

Por ello,

\[
\boxed{
\text{pendiente de la tangente}
=
\text{velocidad instantánea}
}.
\]

Matemáticamente,

\[ v_x=\frac{\dd x}{\dd t}. \]


\subsection{Cómo leer cualitativamente la pendiente}

La forma de la gráfica permite interpretar el movimiento sin realizar
cálculos.

Si la pendiente es positiva,

\[ v_x>0, \]

y la partícula se mueve en el sentido positivo.

Si la pendiente es negativa,

\[ v_x<0, \]

y la partícula se mueve en el sentido negativo.

Si la pendiente es cero,

\[ v_x=0, \]

y la partícula se encuentra instantáneamente en reposo.

Además, cuanto mayor sea el valor absoluto de la pendiente, mayor será la
rapidez.

\begin{idea}
En una gráfica \(x-t\):

\begin{itemize}
    \item pendiente positiva \(\rightarrow\) velocidad positiva;
    \item pendiente negativa \(\rightarrow\) velocidad negativa;
    \item pendiente nula \(\rightarrow\) reposo;
    \item mayor inclinación en valor absoluto \(\rightarrow\) mayor rapidez.
\end{itemize}
\end{idea}


\begin{advertencia}
Una gráfica posición-tiempo no representa necesariamente el dibujo de la
trayectoria.

En movimiento rectilíneo la trayectoria física es una línea recta.

La gráfica \(x-t\) muestra cómo cambia la posición con el tiempo.
\end{advertencia}


\section{Movimiento rectilíneo uniforme}

Hasta ahora hemos descrito movimientos rectilíneos de manera general.

Consideremos ahora un caso particularmente importante: aquel en el cual la
velocidad no cambia.

\begin{definicion}[Movimiento rectilíneo uniforme]
El \textbf{movimiento rectilíneo uniforme}, abreviado MRU, es un movimiento
cuya trayectoria es una línea recta y cuya velocidad permanece constante.
\end{definicion}

Que la velocidad sea constante significa que no cambian:

\begin{itemize}
    \item su magnitud;
    \item su dirección;
    \item su sentido.
\end{itemize}

En consecuencia, una partícula en MRU experimenta desplazamientos iguales
durante intervalos de tiempo iguales.

Por ejemplo, si

\[ v_x=\SI{5}{\meter\per\second}, \]

la partícula experimenta cada segundo un desplazamiento de

\[ +\SI{5}{\meter}. \]

Si

\[ v_x=-\SI{5}{\meter\per\second}, \]

cada segundo experimenta un desplazamiento de

\[ -\SI{5}{\meter}. \]


\subsection{Construcción de la ecuación del MRU}

En MRU la velocidad es constante.

Por tanto, la velocidad media en cualquier intervalo coincide con esa
velocidad constante:

\[ v_x=\frac{\Delta x}{\Delta t}. \]

Multiplicando por \(\Delta t\),

\[ \Delta x=v_x\Delta t. \]

Como

\[ \Delta x=x_f-x_i, \]

tenemos

\[ x_f-x_i=v_x(t_f-t_i). \]

Por tanto,

\[ \boxed{x_f=x_i+v_x(t_f-t_i)}. \]

Esta es la ecuación general de posición para un movimiento rectilíneo con
velocidad constante.

Si elegimos

\[ t_i=0 \]

y escribimos la posición inicial como \(x_0\),

\[ \boxed{x(t)=x_0+v_xt}. \]

\begin{interpretacion}
La expresión

\[ x(t)=x_0+v_xt \]

dice que la posición actual se obtiene sumando:

\[
\text{posición inicial}
+
\text{cambio de posición producido por el movimiento}.
\]

El término

\[ v_xt \]

representa el desplazamiento realizado desde \(t=0\).
\end{interpretacion}


\subsection{¿Qué significa cada término?}

En

\[ x=x_0+v_xt, \]

tenemos:

\begin{itemize}
    \item \(x\): posición en el instante \(t\);
    \item \(x_0\): posición inicial;
    \item \(v_x\): velocidad constante;
    \item \(t\): tiempo transcurrido desde el instante inicial elegido.
\end{itemize}

Cada término posee dimensiones de longitud:

\[
[x]=[x_0]=\si{\meter}
\]

y

\[
[v_xt]
=
\left(\si{\meter\per\second}\right)
(\si{\second})
=
\si{\meter}.
\]

La ecuación es, por tanto, dimensionalmente homogénea.


\begin{advertencia}
No conviene memorizar simplemente

\[ d=vt \]

como ecuación universal del movimiento.

La relación físicamente más completa es

\[ \Delta x=v_x\Delta t. \]

El desplazamiento puede ser positivo o negativo.

La distancia recorrida es una magnitud escalar y debe analizarse
separadamente.
\end{advertencia}


\section{Gráfica posición-tiempo en MRU}

Si

\[ x=x_0+v_xt, \]

la posición depende linealmente del tiempo.

Matemáticamente, esta ecuación tiene la forma

\[ y=b+mx, \]

donde la pendiente corresponde a la velocidad.

En una gráfica \(x-t\):

\[
\boxed{\text{pendiente}=v_x}.
\]

Por tanto, un MRU siempre produce una línea recta en la gráfica
posición-tiempo.


\subsection{Velocidad positiva}

Si

\[ v_x>0, \]

la recta tiene pendiente positiva.

La posición aumenta con el tiempo.

\begin{figure}[H]
    \centering

    \begin{tikzpicture}
        \begin{axis}[
            axis lines=middle,
            xlabel={$t$},
            ylabel={$x$},
            xmin=0,
            xmax=6,
            ymin=0,
            ymax=14,
            width=9cm,
            height=6cm
        ]

            \addplot[
                domain=0:5,
                samples=2,
                thick
            ] {2+2*x};

        \end{axis}
    \end{tikzpicture}

    \caption{Gráfica \(x-t\) para un MRU con velocidad positiva.}
\end{figure}


\subsection{Velocidad negativa}

Si

\[ v_x<0, \]

la recta tiene pendiente negativa.

La posición disminuye con el tiempo.

\begin{figure}[H]
    \centering

    \begin{tikzpicture}
        \begin{axis}[
            axis lines=middle,
            xlabel={$t$},
            ylabel={$x$},
            xmin=0,
            xmax=6,
            ymin=-2,
            ymax=12,
            width=9cm,
            height=6cm
        ]

            \addplot[
                domain=0:5,
                samples=2,
                thick
            ] {10-2*x};

        \end{axis}
    \end{tikzpicture}

    \caption{Gráfica \(x-t\) para un MRU con velocidad negativa.}
\end{figure}


\subsection{Reposo}

Un objeto en reposo tiene

\[ v_x=0. \]

Su posición permanece constante:

\[ x=x_0. \]

La gráfica posición-tiempo es una línea horizontal.

\begin{figure}[H]
    \centering

    \begin{tikzpicture}
        \begin{axis}[
            axis lines=middle,
            xlabel={$t$},
            ylabel={$x$},
            xmin=0,
            xmax=6,
            ymin=0,
            ymax=8,
            width=9cm,
            height=6cm
        ]

            \addplot[
                domain=0:5,
                samples=2,
                thick
            ] {5};

        \end{axis}
    \end{tikzpicture}

    \caption{Gráfica \(x-t\) para una partícula en reposo.}
\end{figure}


\section{Gráfica velocidad-tiempo en MRU}

En MRU la velocidad permanece constante.

Por tanto, una gráfica de velocidad contra tiempo es una línea horizontal.

\begin{figure}[H]
    \centering

    \begin{tikzpicture}
        \begin{axis}[
            axis lines=middle,
            xlabel={$t$},
            ylabel={$v_x$},
            xmin=0,
            xmax=6,
            ymin=0,
            ymax=8,
            width=9cm,
            height=6cm
        ]

            \addplot[
                domain=0:5,
                samples=2,
                thick
            ] {4};

        \end{axis}
    \end{tikzpicture}

    \caption{Gráfica \(v_x-t\) para un MRU con velocidad constante positiva.}
\end{figure}

Durante un intervalo \(\Delta t\), el área rectangular bajo la gráfica es

\[
A=v_x\Delta t.
\]

Pero para MRU,

\[
\Delta x=v_x\Delta t.
\]

Por tanto,

\[ \boxed{\Delta x=A}. \]

\begin{interpretacion}
En una gráfica velocidad-tiempo, el área algebraica comprendida entre la
gráfica y el eje del tiempo representa el desplazamiento.

En MRU esta relación aparece inmediatamente porque el área es simplemente la
de un rectángulo.
\end{interpretacion}

\begin{nota}
Más adelante ampliaremos esta interpretación a movimientos cuya velocidad
no sea constante.
\end{nota}


\section{Ejemplos de movimiento rectilíneo uniforme}

\begin{ejemplo}[Determinar la posición]

Un automóvil se mueve con velocidad constante

\[ v_x=\SI{15}{\meter\per\second}. \]

En \(t=0\) se encuentra en

\[ x_0=\SI{20}{\meter}. \]

Determine su posición después de

\[ t=\SI{8}{\second}. \]

Como la velocidad es constante, utilizamos

\[ x=x_0+v_xt. \]

Sustituyendo,

\[
x
=
\SI{20}{\meter}
+
\left(\SI{15}{\meter\per\second}\right)
\left(\SI{8}{\second}\right).
\]

Entonces,

\[
x
=
\SI{20}{\meter}
+
\SI{120}{\meter}.
\]

Por tanto,

\[ \boxed{x=\SI{140}{\meter}}. \]

\begin{interpretacion}
El automóvil comenzó a veinte metros del origen y realizó un desplazamiento
de ciento veinte metros en el sentido positivo.
\end{interpretacion}
\end{ejemplo}


\begin{ejemplo}[Velocidad negativa]

Una partícula se encuentra inicialmente en

\[ x_0=\SI{30}{\meter} \]

y se mueve con velocidad constante

\[ v_x=-\SI{4.0}{\meter\per\second}. \]

Determine su posición después de

\[ t=\SI{6}{\second}. \]

Aplicamos

\[ x=x_0+v_xt. \]

Entonces,

\[
x
=
\SI{30}{\meter}
+
\left(-\SI{4.0}{\meter\per\second}\right)
\left(\SI{6}{\second}\right).
\]

Por tanto,

\[
x
=
\SI{30}{\meter}
-
\SI{24}{\meter}
=
\SI{6}{\meter}.
\]

Así,

\[ \boxed{x=\SI{6}{\meter}}. \]

La partícula sigue encontrándose en una posición positiva, aunque su
velocidad sea negativa.

\begin{interpretacion}
Este ejemplo muestra nuevamente que el signo de la posición y el signo de la
velocidad describen cosas diferentes.

La partícula está en el lado positivo del origen, pero se mueve en el sentido
negativo.
\end{interpretacion}
\end{ejemplo}


\begin{ejemplo}[Determinar cuándo alcanza una posición]

Un ciclista se mueve en línea recta con velocidad constante

\[ v_x=\SI{6.0}{\meter\per\second}. \]

En \(t=0\) se encuentra en

\[ x_0=-\SI{12}{\meter}. \]

¿En qué instante alcanza la posición

\[ x=\SI{30}{\meter}? \]

Partimos de

\[ x=x_0+v_xt. \]

Despejamos el tiempo:

\[ t=\frac{x-x_0}{v_x}. \]

Sustituimos:

\[
t
=
\frac{
\SI{30}{\meter}-(-\SI{12}{\meter})
}{
\SI{6.0}{\meter\per\second}
}.
\]

Entonces,

\[
t
=
\frac{\SI{42}{\meter}}
{\SI{6.0}{\meter\per\second}}
=
\SI{7.0}{\second}.
\]

Por tanto,

\[ \boxed{t=\SI{7.0}{\second}}. \]
\end{ejemplo}


\section{Encuentro de dos móviles}

La ecuación de posición permite estudiar situaciones con varios móviles.

La idea fundamental es sencilla:

\begin{quote}
dos móviles se encuentran cuando tienen la misma posición en el mismo
instante.
\end{quote}

Por tanto, si sus posiciones son

\[ x_A(t) \]

y

\[ x_B(t), \]

la condición de encuentro es

\[ \boxed{x_A(t)=x_B(t)}. \]


\begin{ejemplo}[Dos móviles que se aproximan]

Dos ciclistas se encuentran inicialmente separados
\(\SI{100}{\meter}\) sobre una carretera recta.

Elegimos el origen en la posición inicial del ciclista A y el sentido
positivo hacia el ciclista B.

El ciclista A se mueve con velocidad constante

\[ v_A=+\SI{6}{\meter\per\second}, \]

mientras que el ciclista B comienza en

\[ x_{B0}=\SI{100}{\meter} \]

y se mueve hacia A con

\[ v_B=-\SI{4}{\meter\per\second}. \]

La posición de A es

\[ x_A=6t. \]

La posición de B es

\[ x_B=100-4t. \]

En el encuentro,

\[ x_A=x_B. \]

Por tanto,

\[
6t=100-4t.
\]

Agrupando,

\[
10t=100.
\]

Así,

\[ t=\SI{10}{\second}. \]

La posición de encuentro es

\[
x_A
=
\left(\SI{6}{\meter\per\second}\right)
\left(\SI{10}{\second}\right)
=
\SI{60}{\meter}.
\]

Por tanto,

\[
\boxed{t_{\text{encuentro}}=\SI{10}{\second}}
\]

y

\[
\boxed{x_{\text{encuentro}}=\SI{60}{\meter}}.
\]

\begin{interpretacion}
No fue necesario memorizar una fórmula especial para ``problemas de
encuentro''.

Bastó escribir correctamente la ecuación de posición de cada móvil y utilizar
la condición física de que ambos ocupan el mismo lugar en el mismo instante.
\end{interpretacion}
\end{ejemplo}


\section{Procedimiento para resolver problemas de MRU}

\begin{notapedagogica}
Para evitar que los problemas de movimiento se conviertan en sustituciones
automáticas, puede enseñarse el siguiente procedimiento:

\begin{enumerate}

    \item Dibujar un esquema sencillo de la situación.

    \item Elegir el origen.

    \item Elegir el sentido positivo del eje.

    \item Escribir las posiciones iniciales con sus signos.

    \item Asignar el signo correcto a cada velocidad.

    \item Identificar qué magnitud se busca.

    \item Escribir la ecuación de posición:

    \[
    x=x_0+v_xt.
    \]

    \item Sustituir los datos con unidades coherentes.

    \item Resolver algebraicamente.

    \item Interpretar el signo y el valor obtenido.

\end{enumerate}

La elección del sistema de referencia debe hacerse antes de introducir los
signos en las ecuaciones.
\end{notapedagogica}


\section{Cuándo puede utilizarse el modelo de MRU}

El movimiento rectilíneo uniforme es un modelo idealizado.

\begin{condicion}
La ecuación

\[ x=x_0+v_xt \]

puede utilizarse durante un intervalo en el que:

\begin{itemize}
    \item la trayectoria sea recta;
    \item la velocidad pueda considerarse constante;
    \item se mantenga el mismo sistema de referencia;
    \item las posiciones y tiempos se midan respecto de los orígenes
    establecidos.
\end{itemize}
\end{condicion}

Muchos movimientos reales no mantienen exactamente una velocidad constante.

Sin embargo, el MRU puede proporcionar una descripción adecuada durante
intervalos en los que las variaciones de velocidad sean despreciables para
el propósito del análisis.


\section{Errores y confusiones frecuentes}

\begin{errorfrecuente}
\textbf{Confundir velocidad media con rapidez media.}

La velocidad media utiliza desplazamiento:

\[ \bar{v}_x=\frac{\Delta x}{\Delta t}. \]

La rapidez media utiliza distancia:

\[ \bar{v}_{\text{rap}}=\frac{s}{\Delta t}. \]
\end{errorfrecuente}


\begin{errorfrecuente}
\textbf{Pensar que una velocidad negativa significa una rapidez negativa.}

La rapidez nunca es negativa.

El signo de una componente de velocidad indica el sentido del movimiento.
\end{errorfrecuente}


\begin{errorfrecuente}
\textbf{Confundir posición negativa con velocidad negativa.}

Un objeto puede encontrarse en \(x<0\) y moverse con \(v_x>0\).

También puede encontrarse en \(x>0\) y moverse con \(v_x<0\).
\end{errorfrecuente}


\begin{errorfrecuente}
\textbf{Pensar que la velocidad media describe necesariamente la velocidad
en todos los instantes.}

Esto solo ocurre en el caso particular en que la velocidad permanece
constante.
\end{errorfrecuente}


\begin{errorfrecuente}
\textbf{Calcular rapidez media utilizando el desplazamiento.}

Si el móvil cambia de sentido, debe calcularse primero la distancia total
recorrida.
\end{errorfrecuente}


\begin{errorfrecuente}
\textbf{Interpretar la altura de una gráfica \(x-t\) como velocidad.}

En una gráfica posición-tiempo, la coordenada vertical representa posición.

La velocidad está relacionada con la \textbf{pendiente}.
\end{errorfrecuente}


\begin{errorfrecuente}
\textbf{Pensar que una recta horizontal en una gráfica \(x-t\) representa
movimiento uniforme.}

Una recta horizontal significa

\[ x=\text{constante}, \]

por tanto,

\[ v_x=0. \]

El cuerpo está en reposo.
\end{errorfrecuente}


\begin{errorfrecuente}
\textbf{Usar \(d=vt\) sin definir dirección ni signos.}

En problemas unidimensionales es preferible trabajar con

\[ \Delta x=v_x\Delta t \]

o con

\[ x=x_0+v_xt. \]

Estas expresiones conservan explícitamente la información sobre el sentido
del movimiento.
\end{errorfrecuente}


\section{Secuencia conceptual para explicar el tema}

\begin{notapedagogica}
Una secuencia adecuada para dictar esta clase es:

\begin{enumerate}

    \item Recordar brevemente posición, desplazamiento, distancia e intervalo
    de tiempo.

    \item Plantear la pregunta:

    \begin{quote}
    ¿Cómo podemos expresar qué tan rápidamente cambia la posición?
    \end{quote}

    \item Introducir la velocidad media como desplazamiento dividido entre
    tiempo.

    \item Analizar inmediatamente el signo de la velocidad media.

    \item Resolver un ejemplo con velocidad media positiva y otro con
    velocidad media negativa.

    \item Introducir rapidez media y compararla con velocidad media.

    \item Resolver un ejemplo con cambio de sentido para mostrar que
    rapidez media y magnitud de la velocidad media no son iguales en
    general.

    \item Preguntar qué información falta si solamente conocemos la velocidad
    media durante un intervalo largo.

    \item Introducir intuitivamente la velocidad instantánea mediante
    intervalos cada vez más pequeños.

    \item Presentar formalmente

    \[
    v_x=\frac{\dd x}{\dd t}.
    \]

    \item Relacionar rapidez instantánea con

    \[
    v=\abs{v_x}.
    \]

    \item Introducir la gráfica posición-tiempo.

    \item Explicar que la pendiente de una secante representa velocidad
    media.

    \item Explicar que la pendiente de una tangente representa velocidad
    instantánea.

    \item Interpretar gráficamente velocidades positivas, negativas y nulas.

    \item Introducir el MRU como el caso especial de velocidad constante.

    \item Construir, en lugar de memorizar, la ecuación

    \[
    x=x_0+v_xt.
    \]

    \item Interpretar cada término físicamente.

    \item Dibujar las gráficas \(x-t\) y \(v-t\) del MRU.

    \item Resolver ejemplos con velocidad positiva y negativa.

    \item Resolver un problema de encuentro utilizando dos ecuaciones de
    posición.

    \item Finalizar señalando que en la siguiente etapa estudiaremos qué
    ocurre cuando la velocidad deja de ser constante.

\end{enumerate}
\end{notapedagogica}


\section{Preguntas para comprobar comprensión}

Antes de continuar con aceleración pueden plantearse preguntas como:

\begin{enumerate}

    \item ¿Cuál es la diferencia entre desplazamiento y distancia recorrida?

    \item ¿Cuál es la diferencia entre velocidad media y rapidez media?

    \item ¿Puede una velocidad media ser negativa?

    \item ¿Puede una rapidez media ser negativa?

    \item ¿Puede un objeto tener rapidez media distinta de cero y velocidad
    media igual a cero?

    \item ¿Qué significa físicamente una velocidad negativa?

    \item Si \(x<0\), ¿podemos concluir que \(v_x<0\)?

    \item ¿Qué representa la pendiente de una recta secante en una gráfica
    \(x-t\)?

    \item ¿Qué representa la pendiente de una recta tangente?

    \item ¿Qué significa una pendiente horizontal en una gráfica \(x-t\)?

    \item ¿Qué movimiento representa una línea recta inclinada en una gráfica
    \(x-t\)?

    \item Si dos rectas \(x-t\) tienen pendientes positivas, ¿cómo podemos
    saber qué móvil tiene mayor rapidez?

    \item ¿Qué significa que un movimiento sea uniforme?

    \item Si un objeto tiene velocidad constante negativa, ¿su posición
    aumenta o disminuye con el tiempo?

    \item ¿Por qué en un MRU la gráfica \(x-t\) debe ser una línea recta?

    \item ¿Qué condición matemática permite encontrar el instante en que dos
    móviles se encuentran?

\end{enumerate}


\section{Conexión con el tema siguiente}

El movimiento rectilíneo uniforme constituye un caso especial:

\[ v_x=\text{constante}. \]

Sin embargo, la mayoría de los movimientos reales presentan cambios de
velocidad.

Un automóvil puede arrancar, aumentar su rapidez, frenar o cambiar el sentido
de su movimiento.

La pregunta natural es entonces:

\begin{quote}
¿Cómo describimos qué tan rápidamente cambia la velocidad?
\end{quote}

La respuesta conducirá al concepto de \textbf{aceleración}.

A partir de allí estudiaremos el movimiento rectilíneo con aceleración
constante.


\begin{resumen}

\begin{itemize}

    \item En movimiento rectilíneo la posición puede describirse mediante una
    sola coordenada:

    \[
    x=x(t).
    \]

    \item La velocidad media es

    \[
    \boxed{\bar{v}_x=\frac{\Delta x}{\Delta t}}.
    \]

    \item La velocidad media depende del desplazamiento y puede ser positiva,
    negativa o cero.

    \item La rapidez media es

    \[
    \boxed{\bar{v}_{\text{rap}}=\frac{s}{\Delta t}}.
    \]

    \item La rapidez media depende de la distancia total recorrida y nunca es
    negativa.

    \item En general,

    \[
    \bar{v}_{\text{rap}}\neq\abs{\bar{v}_x}.
    \]

    \item La velocidad instantánea se define como

    \[
    \boxed{
    v_x
    =
    \lim_{\Delta t\to0}
    \frac{\Delta x}{\Delta t}
    =
    \frac{\dd x}{\dd t}
    }.
    \]

    \item La rapidez instantánea es la magnitud de la velocidad:

    \[
    \boxed{v=\abs{v_x}}.
    \]

    \item En una gráfica posición-tiempo, la pendiente de una secante
    representa la velocidad media.

    \item La pendiente de la tangente representa la velocidad instantánea.

    \item Una pendiente positiva indica velocidad positiva.

    \item Una pendiente negativa indica velocidad negativa.

    \item Una pendiente nula indica reposo instantáneo.

    \item El movimiento rectilíneo uniforme tiene trayectoria recta y
    velocidad constante.

    \item En MRU,

    \[
    \boxed{\Delta x=v_x\Delta t}.
    \]

    \item La ecuación general de posición es

    \[
    \boxed{x_f=x_i+v_x(t_f-t_i)}.
    \]

    \item Si \(t_i=0\),

    \[
    \boxed{x(t)=x_0+v_xt}.
    \]

    \item La gráfica \(x-t\) de un MRU es una línea recta cuya pendiente es
    la velocidad.

    \item La gráfica \(v_x-t\) de un MRU es una línea horizontal.

    \item En un problema con dos móviles, el encuentro ocurre cuando

    \[
    \boxed{x_A(t)=x_B(t)}.
    \]

\end{itemize}

\end{resumen}